% savepos-warichu.tex -- acceptance gate: savepos-driven machinery
% (interlinear reference marks + 割注 warichu) under d-tate, with the
% logical-frame line-extent computation checked against the class's
% own geometry instead of the raw PDF X coordinate.
%
% The aux frame-id round-trip: \pdfsavepos positions are recorded per
% \jlreq@reference@count id in the aux; pass 1 has no origins yet
% (labels default to 0), so the geometric assertions are made on the
% FINAL pass (the runner runs 2 passes and asserts on pass 2's log),
% proving pass 2 consumes pass 1's origins.
%
% Assertions (log markers, asserted on pass 2 by the runner):
%   MARK-POS  : an interlinear reference mark placed mid-line keeps the
%               pair glue of the surrounding chars (same natural-width
%               contract as the yoko interlinear fixture, in tate).
%   REST-WIDTH: the warichu's internal savepos measurement (taken by
%               the class at the warichu start) reports the rest of the
%               logical line within tolerance: line 20zw, run of 5
%               kanji + opening bracket before the warichu cell.
%   WARICHU-OK: the warichu typeset (2 lines) and the paragraph ships.
\documentclass[pdflatex,tate,paper=b5,line_length=20zw]{jlreq}
\jlreqsetup{reference_mark=interlinear}
\makeatletter
\begin{document}
% ---- MARK-POS: 合印を挟んでも前後のペアグルーが不変（横組版と同一契約）
\newbox\mk@boxa \newbox\mk@boxb
\setbox\mk@boxa=\hbox{あ、\footnote{注1}あ}%
\setbox\mk@boxb=\hbox{あ、あ}%
\@tempdima=\dimexpr\wd\mk@boxa-\wd\mk@boxb\relax
\ifdim\@tempdima<\z@ \@tempdima=-\@tempdima\fi
\ifdim\@tempdima<0.5pt
  \typeout{MARK-POS: pair glue intact (\the\@tempdima)}%
\else
  \typeout{MARK-POS-FAIL: delta=\the\@tempdima}%
\fi

% ---- REST-WIDTH / WARICHU-OK: 割注（コマンド形式）。
% 20zw の行に全角5文字を置いてから割注を開始。割注内部で
% \jlreq@savepos@for@rest@linewidth が走り、
% \jlreq@current@linewidth / \jlreq@rest@linewidth に
% 論理行方向（=内枠 X 方向）の値が入る。
あいうえお\warichu{割注一行目　割注二行目}残りの本文。
\typeout{REST-WIDTH: current=\the\jlreq@current@linewidth,
  rest=\the\jlreq@rest@linewidth, linewidth=\the\linewidth,
  zw=\the\jlreq@zw}%
続きの文章。ここで検査は終わる。\footnote{脚注も縦組みで出る。}
\end{document}
